% =====================================================================
% RASAD: tables for the paper (IEEEtran, two-column)
% Preamble needs:  \usepackage{booktabs}
%
% Every number is either reported in the paper or derived from the
% Fig. 6 confusion matrix with compute_metrics.py. Results already shown
% in figures (confusion matrix, modality bars, ROC AUCs) are not repeated.
% =====================================================================


% ---------------------------------------------------------------------
% Table II: dataset (place in Sec. III, after the ADReSSo21 paragraph)
% ---------------------------------------------------------------------
\begin{table}[!t]
\centering
\caption{ADReSSo21 data used in this study}
\label{tab:dataset}
\begin{tabular}{lccc}
\toprule
Split & AD & CN & Total \\
\midrule
Training (5-fold CV) & 87 & 79 & 166 \\
Test (held out)      & 35 & 36 & 71  \\
\midrule
Total                & 122 & 115 & 237 \\
\bottomrule
\end{tabular}
\end{table}


% ---------------------------------------------------------------------
% Table III: test-set metrics with uncertainty (Sec. IV-B)
% Derived from the Fig. 6 confusion matrix. 95% CIs are Wilson score
% intervals; they quantify the uncertainty stated in the Conclusion.
% ---------------------------------------------------------------------
\begin{table}[!t]
\centering
\caption{Test-set performance of the proposed ensemble ($n=71$; AD is the positive class)}
\label{tab:metrics}
\begin{tabular}{lcc}
\toprule
Metric & Value & 95\% CI \\
\midrule
Accuracy                  & 84.5\% & 74.3--91.1\% \\
Sensitivity (AD recall)   & 82.9\% & 67.3--91.9\% \\
Specificity (CN recall)   & 86.1\% & 71.3--93.9\% \\
Precision (PPV)           & 85.3\% & 69.9--93.6\% \\
NPV                       & 83.8\% & 68.9--92.3\% \\
\midrule
F1-score (AD)             & 0.841 & -- \\
Macro F1-score            & 0.845 & -- \\
Matthews corr.\ coef.     & 0.690 & -- \\
ROC AUC                   & 0.89  & -- \\
Average precision         & 0.89  & -- \\
\bottomrule
\end{tabular}
\end{table}


% ---------------------------------------------------------------------
% Table IV: same-benchmark comparison (Sec. IV, or end of Sec. II-G)
% Only systems evaluated on the ADReSSo21 test set, so accuracies are
% directly comparable (unlike Table I, which mixes corpora).
% Verify each figure against the cited paper before submission.
% ---------------------------------------------------------------------
\begin{table}[!t]
\centering
\caption{Comparison with published systems on the ADReSSo21 test set}
\label{tab:adresso}
\begin{tabular}{lllc}
\toprule
System & Modalities & Inference & Acc.\ (\%) \\
\midrule
Luz et al.~\cite{b35} (baseline) & A + L & Offline & 78.87 \\
Zhu et al.~\cite{b20} (WavBERT)  & A + L & Offline & 83.10 \\
Rohanian et al.~\cite{b15}       & A + L & Offline & 84.00 \\
Pan et al.~\cite{b8}             & L     & Offline & 84.51 \\
Syed et al.~\cite{b18}           & L     & Offline & 84.51 \\
\midrule
\textbf{RASAD (ours)}            & A + L & Real-time & \textbf{84.51} \\
\bottomrule
\multicolumn{4}{l}{\footnotesize A: acoustic; L: linguistic.}
\end{tabular}
\end{table}
% Adjust \cite keys (b8, b15, ...) to match your .bib / thebibliography labels.


% ---------------------------------------------------------------------
% Table V: expected screening performance at realistic prevalence
% (Sec. IV-B or Discussion). ADReSSo21 is balanced (49% AD in the test
% set); screened populations are not. PPV and NPV follow from the test
% sensitivity (82.9%) and specificity (86.1%) by Bayes' rule; counts are
% per 1,000 people screened. Computed by compute_metrics.py.
% ---------------------------------------------------------------------
\begin{table}[!t]
\centering
\caption{Projected screening performance of RASAD at different AD prevalences (per 1{,}000 people screened)}
\label{tab:screening}
\begin{tabular}{cccccc}
\toprule
Prevalence & PPV (\%) & NPV (\%) & TP & FP & FN \\
\midrule
5\%  & 23.9 & 99.0 & 41  & 132 & 9  \\
10\% & 39.9 & 97.8 & 83  & 125 & 17 \\
20\% & 59.9 & 95.3 & 166 & 111 & 34 \\
30\% & 71.9 & 92.1 & 249 & 97  & 51 \\
50\% & 85.6 & 83.4 & 414 & 69  & 86 \\
\bottomrule
\multicolumn{6}{l}{\footnotesize Positive likelihood ratio 5.97; negative likelihood ratio 0.20.}
\end{tabular}
\end{table}
